% Part: sets-functions-relations
% Chapter: arithmetization
% Section: cauchy

\documentclass[../../../include/open-logic-section]{subfiles}

\begin{document}

\olfileid{sfr}{arith}{cauchy}

\olsection{পরিশিষ্ট: কোশি অনুক্রম হিসেবে বাস্তব সংখ্যা}

\olref[cuts]{sec}-এ আমরা বাস্তব সংখ্যাকে ডেডেকিন্ড কর্তন হিসেবে নির্মাণ করেছি। এই অংশে একটি বিকল্প নির্মাণ ব্যাখ্যা করি। এটি কোশির (যাকে এখন আমরা বলি) কোশি অনুক্রমের সংজ্ঞার উপর গড়ে উঠেছে; তবে এই সংজ্ঞা ব্যবহার করে বাস্তব সংখ্যা \emph{নির্মাণ} করার কৃতিত্ব ঊনবিংশ শতাব্দীর অন্য লেখকদের, বিশেষত ভাইয়ারস্ট্রাস, হাইনে, মেরি ও কান্টরের। (সুন্দর একটি ইতিহাসের জন্য \citeauthor{OConnorRobertson:RN}
\citeyear{OConnorRobertson:RN} দেখো।)

ঊনবিংশ শতাব্দীতে যাওয়ার আগে সাইমন স্টেভিনকে (1548--1620) বিবেচনা করা দরকার। সংক্ষেপে, স্টেভিন উপলব্ধি করেছিলেন যে প্রত্যেক বাস্তব সংখ্যাকে তার দশমিক বিস্তারের মাধ্যমে ভাবা যায়। তাই $\sqrt{2}$-এর মতো অমূলদ সংখ্যারও সুন্দর একটি দশমিক বিস্তার আছে, যার শুরু:
\[
	1.41421356237\ldots
\]
সেটতত্ত্বে দশমিক বিস্তারের মডেল তৈরি করা খুব সহজ: এগুলিকে স্রেফ অপেক্ষক $d \colon \Nat \to \Nat$ হিসেবে বিবেচনা করো, যেখানে $d(n)$ হলো আমাদের বিবেচ্য $n$-তম দশমিক স্থান। তখন দশমিক বিন্দুর আগে থাকা বাস্তব সংখ্যার অংশটি (এখানে শুধু $1$) সামলাতে সংজ্ঞায় একটু পরিবর্তন দরকার হবে। যেমন $0.999\ldots
= 1$ নিশ্চিত করতে আরও একটি পরিবর্তন (একটি তুল্যতা সম্পর্ক) লাগবে। তবে স্টেভিনের পদ্ধতিতে সেটতত্ত্বের মধ্যে বাস্তব সংখ্যার সম্পূর্ণ নিখুঁত নির্মাণ দেওয়া কঠিন নয়।

স্টেভিন আমাদের মূল বিষয় নন। (স্টেভিন সম্পর্কে আরও জানতে \citealt{KatzKatz2012} দেখো।) তবে ঘনিষ্ঠভাবে সম্পর্কিত একটি চিন্তা এখানে। $\sqrt{2}$-এর দশমিক বিস্তারটি সরাসরি না ধরে, ক্রমশ আরও নিখুঁত বিস্তারগুলি নিয়ে $\sqrt{2}$-এর ক্রমশ আরও যথাযথ মূলদ আসন্নমানের একটি \emph{অনুক্রম} বিবেচনা করতে পারি:
\[
1, 1.4, 1.414, 1.4142, 1.41421,\ldots
\]
বাস্তব সংখ্যাকে ``ক্রমশ আরও ভালো আসন্নমান'' দিয়ে বিবেচনা করার ধারণাটি আমাদের আরও এক ধারার অন্তর্দৃষ্টির ভিত্তি দেয় (যেমন $\Int$ থেকে $\Rat$, অথবা $\Nat$ থেকে $\Int$ নির্মাণের সময় ব্যবহৃত উপলব্ধিগুলি)। মৌলিক অন্তর্দৃষ্টিগুলি হলো:
\begin{enumerate}
	\item প্রত্যেক বাস্তব সংখ্যাকে একটি (সম্ভবত অসীম) দশমিক বিস্তার হিসেবে লেখা যায়।
	\item একটি (সম্ভবত অসীম) দশমিক বিস্তারে সংকেতায়িত তথ্য মূলদ সংখ্যার একটি অনুক্রম দিয়েও একইভাবে সংকেতায়িত করা যায়।
	\item মূলদ সংখ্যার একটি অনুক্রমকে $\Nat$ থেকে $\Rat$-এ একটি অপেক্ষক হিসেবে ভাবা যায়; শুধু $f(n)$-কে অনুক্রমের $n$-তম মূলদ সংখ্যাটি নাও।
\end{enumerate}
অবশ্য $\Nat$ থেকে $\Rat$-এ যেকোনো অপেক্ষকই আমাদের একটি বাস্তব সংখ্যা দেবে না। যেমন, এই অপেক্ষকটি বিবেচনা করো:
\[
	f(n) =\begin{cases}
			1 & \text{যদি }n\text{ বিজোড় হয়}\\
			0 &\text{যদি }n\text{ জোড় হয়}
		\end{cases}
\]
এখানে মূল সমস্যা হলো, অনুক্রম $0,1,0,1,0,1,0,\ldots$ কোনো বাস্তব সংখ্যার দিকে ``ঘনিয়ে যায়'' বলে মনে হয় না। তাই কোনো বাস্তব সংখ্যার দিকে সত্যিই ঘনিয়ে যায় এমন অনুক্রম বিবেচনা করতে আমাদের মনোযোগ সেইসব অনুক্রমে সীমাবদ্ধ করতে হবে, যাদের কোনো \emph{সীমা} আছে।

\olref[his][set][limits]{sec}-এ আমরা সীমার ধারণা ইতিমধ্যেই দেখেছি। কিন্তু সেখানে যে সংজ্ঞা ব্যবহার করেছি, ঠিক \emph{সেই} সংজ্ঞাটি এখানে ব্যবহার করা যাবে না। সেখানে ``$(\forall \epsilon>0)$'' রাশিটির মধ্যে প্রচ্ছন্নভাবে বাস্তব সংখ্যার উপর পরিমাণায়ন ছিল; আর আমরা বাস্তব সংখ্যার উপর অপেক্ষকের সীমা বিবেচনা করছিলাম। তাই ওই সংজ্ঞা প্রয়োগ করলে বাস্তব সংখ্যাগুলিকে আগে থেকেই ধরে নেওয়া হতো, অথচ ঠিক সেগুলিই আমরা \emph{নির্মাণ} করতে চাইছি। সৌভাগ্যক্রমে, সীমার একটি ঘনিষ্ঠ সম্পর্কিত ধারণা নিয়ে কাজ করা যায়।
\begin{defn}\ollabel{def:CauchySequence}
  একটি অপেক্ষক $f: \Nat \to \Rat$ একটি \emph{কোশি অনুক্রম} হয় যদি এবং কেবল যদি যেকোনো ধনাত্মক $\epsilon \in \Rat$-এর জন্য $(\exists \ell \in
  \Nat)(\forall m, n > \ell)|f(m) - f(n)| < \epsilon$ হয়।
\end{defn}

সীমার সাধারণ ধারণাটি আগের মতোই: নির্দিষ্ট মাত্রার নিখুঁততা (যা~$\epsilon$ দিয়ে মাপা) চাইলে দেখার মতো একটি ``অঞ্চল'' আছে (~$\ell$-এর চেয়ে বড় যেকোনো ইনপুট)। আর আমাদের অনুক্রম $1$, $1.4$, $1.414$, $1.4142$,
$1.41421$\ldots-এর যে একটি সীমা আছে, তা সহজেই দেখা যায়: $\nicefrac{1}{10^{n}}$-এর চেয়ে কম ত্রুটিতে $\sqrt{2}$-এর আসন্নমান চাইলে $n$-তম পদের পরের যেকোনো পদ দেখো।

তাহলে স্পষ্ট চিন্তাটি হলো, যেকোনো কোশি অনুক্রমই একটি বাস্তব সংখ্যা \emph{হয়} বলা। কিন্তু $\Int$ ও $\Rat$-এর নির্মাণের মতো এটিও অতিসরল: যেকোনো একটি বাস্তব সংখ্যাকে একাধিক ভিন্ন কোশি অনুক্রম নির্দেশ করে। এটি দেখার সহজ উপায় নিচে। একটি কোশি অনুক্রম~$f$ দেওয়া থাকলে $g$-কে~$f$-এর ঠিক একই অপেক্ষক হিসেবে সংজ্ঞায়িত করো, শুধু $g(0)\neq
f(0)$। প্রথম সংখ্যার পর সর্বত্র অনুক্রমদুটি মিলে যায় বলে, \olref{def:CauchySequence}-এ ব্যবহৃত অর্থে তাদের একই সীমা আছে—এ কথা আমরা (শেষ পর্যন্ত) বলতে চাইব; তাই তাদের একই বাস্তব সংখ্যাকে ``সংজ্ঞায়িত'' করে বলে ভাবা উচিত। অতএব কোশি অনুক্রমদুটিকে আসলে একই বাস্তব সংখ্যা হিসেবে ভাবা উচিত।

ফলে কোশি অনুক্রমগুলির উপর আবার একটি তুল্যতা সম্পর্ক সংজ্ঞায়িত করতে এবং বাস্তব সংখ্যাগুলিকে তুল্যতা শ্রেণির সঙ্গে অভিন্ন করতে হবে। % BN-SRC-013 TARGET-CORRECTION-BEGIN
\par\smallskip\noindent\textnormal{\small\textbf{উৎস-সংশোধন (BN-SRC-013):} হিমায়িত ইংরেজি উৎসের এই বাক্যে বাস্তব সংখ্যাকে ``তুল্যতা সম্পর্ক''-এর সঙ্গে অভিন্ন করার কথা আছে। পরের সংজ্ঞা ও আগের নির্মাণগুলির মতো এখানে অভিপ্রেত বস্তু হলো ওই সম্পর্কের \emph{তুল্যতা শ্রেণি}; বাংলা পাঠে সেটিই বলা হয়েছে।} % BN-SRC-013 TARGET-CORRECTION-END
প্রথমে সীমায় $0$-র দিকে যায় এমন অপেক্ষকের ধারণা দরকার। যেকোনো অপেক্ষক $h : \Nat \to \Rat$-এর জন্য বলি, \emph{$h$
সীমায় $0$-র দিকে যায়} যদি এবং কেবল যদি যেকোনো ধনাত্মক $\epsilon \in \Rat$-এর জন্য $(\exists \ell \in \Nat)(\forall n > \ell)|h(n)| <
\epsilon$ হয়।\footnote{এটি \olref[his][set][limits]{sec}-এর $\lim_{x
\mathord{\rightarrow}\infty}f(x) = 0$ সংজ্ঞার সঙ্গে তুলনা করো।} আরও, $f$ ও $g$ যদি $\Nat \to \Rat$ অপেক্ষক হয়, তবে ধরি $(f-g)(n) = f(n) - g(n)$। এবার সংজ্ঞায়িত করি:
\[
	f \Realequiv g \text{ যদি এবং কেবল যদি $(f-g)$ সীমায় $0$-র দিকে যায়।}.
\]
$\Realequiv$ একটি তুল্যতা সম্পর্ক কি না যাচাই করতে হবে; এবং সেটি তা-ই। তাহলে চাইলে বাস্তব সংখ্যাগুলিকে $\Nat \to
\Rat$ সকল কোশি অনুক্রমের $\Realequiv$-এর অধীনে তুল্যতা শ্রেণি হিসেবে সংজ্ঞায়িত করতে পারি।

\begin{prob}
প্রত্যেক $n$-এর জন্য ধরি $f(n) = 0$। ধরি $g(n) = \frac{1}{(n+1)^2}$। দেখাও যে উভয়েই কোশি অনুক্রম এবং উভয় অপেক্ষকের সীমাই $0$, ফলে $f \sim_\Real g$-ও হয়।
\end{prob}

এটি করার পরে, যথারীতি যার !!a{element} $f$ সেই তুল্যতা শ্রেণির জন্য $\equivrep{f}{\Realequiv}$ লিখব। তবে পাঠযোগ্যতা বজায় রাখতে পরের আলোচনায় সাবস্ক্রিপ্টটি বাদ দিয়ে শুধু $\equivrep{f}{}$ লিখব। আরও নির্ধারণ করি যে প্রত্যেক $q \in \Rat$-এর জন্য $q_{\Real} = \equivrep{c_{q}}{}$, যেখানে $c_{q}$ হলো ধ্রুব অপেক্ষক $c_q(n) = q$, সকল $n \in \Nat$-এর জন্য। এরপর বাস্তব সংখ্যার উপর মৌলিক সম্পর্ক ও ক্রিয়া সংজ্ঞায়িত করি, যেমন:
\begin{align*}
	\equivrep{f}{} + \equivrep{g}{} &= 	\equivrep{(f + g)}{} \\
	\equivrep{f}{} \times 	\equivrep{g}{} &= \equivrep{(f \times g)}{}
\end{align*}
যেখানে $(f + g)(n) = f(n) + g(n)$ এবং $(f \times g)(n) = f(n) \times
g(n)$। অবশ্য $f$ ও $g$ কোশি অনুক্রম হলে $(f + g)$,
$(f-g)$ ও $(f\times g)$-এর প্রত্যেকটিও কোশি অনুক্রম কি না যাচাই করতে হবে; সেগুলি তা-ই, এবং কাজটি তোমার জন্য রেখে দিই।

সবশেষে, ক্রমের একটি ধারণা সংজ্ঞায়িত করি। বলি $\equivrep{f}{}$
\emph{ধনাত্মক} যদি এবং কেবল যদি $\equivrep{f}{}\neq 0_\Real$ এবং $(\exists
\ell \in \Nat)(\forall n > \ell)0 < f(n)$ উভয়ই হয়। % BN-SRC-014 TARGET-CORRECTION-BEGIN
\par\smallskip\noindent\textnormal{\small\textbf{উৎস-সংশোধন (BN-SRC-014):} হিমায়িত ইংরেজি উৎসে বাস্তব তুল্যতা শ্রেণি $\equivrep{f}{}$-কে $0_\Rat$-এর সঙ্গে তুলনা করা হয়েছে। এই নির্মাণে সঠিক একই-প্রকারের শূন্য হলো $0_\Real$; বাংলা সূত্রে সাবস্ক্রিপ্টটি ঠিক করা হয়েছে।} % BN-SRC-014 TARGET-CORRECTION-END
তারপর বলি $\equivrep{f}{} <
\equivrep{g}{}$ যদি এবং কেবল যদি $\equivrep{(g - f)}{}$ ধনাত্মক। এটি সুসংজ্ঞায়িত কি না (অর্থাৎ তুল্যতা শ্রেণি থেকে বেছে নেওয়া ``প্রতিনিধি'' অপেক্ষকের উপর নির্ভর করে না কি না) যাচাই করতে হবে।
%\begin{align*}
%	\equivrep{f}{} \leq \equivrep{g}{} \text{ যদি এবং কেবল যদি }&\equivrep{f}{} = \equivrep{g}{} \text{ অথবা }\\
%	&(\exists \ell \in \Nat)(\forall n \in \Nat)(\ell < n \lif f(n) \leq g(n))
%\end{align*}
এটি করা হয়ে গেলে, এগুলি যে সঠিক বীজগাণিতিক ধর্ম দেয় তা বেশ সহজেই দেখানো যায়; অর্থাৎ:
\begin{thm}\ollabel{thm:cauchyorderedfield}
কোশি অনুক্রমগুলির তুল্যতা শ্রেণিগুলি একটি ক্রমিত ক্ষেত্র গঠন করে। % BN-SRC-015 TARGET-CORRECTION-BEGIN
\par\smallskip\noindent\textnormal{\small\textbf{উৎস-ব্যাখ্যা (BN-SRC-015):} হিমায়িত ইংরেজি উপপাদ্য ও অনুশীলনে সংক্ষেপে ``কোশি অনুক্রমগুলি'' বলা হয়েছে। সংজ্ঞায়িত বাস্তব সত্তা ও ক্রিয়াগুলি আসলে $\Realequiv$-এর অধীনে তাদের তুল্যতা শ্রেণির উপর; বাংলা পাঠে সেই প্রকারটি স্পষ্ট করা হয়েছে।} % BN-SRC-015 TARGET-CORRECTION-END
\end{thm}

\begin{proof}
অনুশীলন।
\end{proof}

\begin{prob}
প্রমাণ করো যে কোশি অনুক্রমগুলির তুল্যতা শ্রেণিগুলি একটি ক্রমিত ক্ষেত্র গঠন করে।
\end{prob}

এইভাবে নির্মিত বাস্তব সংখ্যার পূর্ণতা ধর্ম আছে—এটি প্রমাণ করা কঠিনতর, তাই প্রমাণটি দেব।

\begin{thm}
কোশি অনুক্রমগুলির তুল্যতা শ্রেণির প্রত্যেক অশূন্য ঊর্ধ্বসীমাযুক্ত সেটের একটি লঘিষ্ঠ ঊর্ধ্বসীমা আছে।
\end{thm}

\begin{proof}[প্রমাণের রূপরেখা] ধরি, $S$ কোশি অনুক্রমগুলির যেকোনো অশূন্য ঊর্ধ্বসীমাযুক্ত সেট। তাহলে এমন কোনো $p \in \Rat$ আছে যাতে $p_{\Real}$ হলো $S$-এর একটি ঊর্ধ্বসীমা। ধরি $r \in S$; তাহলে এমন কোনো $q \in \Rat$ আছে যাতে $q_{\Real} < r$। তাই $S$-এর কোনো লঘিষ্ঠ ঊর্ধ্বসীমা থাকলে সেটি $q_\Real$ ও $p_\Real$-এর মধ্যে (প্রান্তদুটিসহ) থাকবে।

নিচ ও উপর থেকে একই সঙ্গে কাছে গিয়ে আমরা লঘিষ্ঠ ঊর্ধ্বসীমাটির দিকে ঘনিয়ে যাব। বিশেষভাবে, দুটি অপেক্ষক $f, g \colon
\Nat \to \Rat$ সংজ্ঞায়িত করি; উদ্দেশ্য হলো $f$ উপর থেকে এবং $g$ নিচ থেকে লঘিষ্ঠ ঊর্ধ্বসীমাটির দিকে ঘনিয়ে যাবে। শুরুতে সংজ্ঞায়িত করি:
\begin{align*}
	f(0) &= p \\
	g(0) &= q
\end{align*}
তারপর $a_n = \frac{f(n) + g(n)}{2}$ হলে ধরি:\footnote{এটি একটি পূর্ববর্তী মানের সাহায্যে ধাপে ধাপে সংজ্ঞা। কিন্তু এমন সংজ্ঞা গ্রহণযোগ্য ভাবার কোনো কারণ আমরা \emph{এখনও} দিইনি।}
\begin{align*}
	f(n+1) &=
	\begin{cases}
		a_n &\text{যদি }(\forall h \in S)\equivrep{h}{} \leq (a_n)_\Real\\
		f(n)&\text{অন্যথায়}
	\end{cases}\\
	g(n+1) &=
	\begin{cases}
		a_n &\text{যদি }(\exists h \in S)\equivrep{h}{} \geq (a_n)_\Real\\
		g(n) &\text{অন্যথায়}
	\end{cases}
\end{align*}
$f$ ও $g$ উভয়েই কোশি অনুক্রম। (এটি বেশ সহজেই যাচাই করা যায়, তবে কাজটি অনুশীলন হিসেবে রেখে দিই।) লক্ষ করো, অপেক্ষক $(f-g)$ সীমায় $0$-র দিকে যায়, কারণ প্রত্যেক ধাপে $f$ ও $g$-এর ব্যবধান আগের ব্যবধানের অর্ধেকের বেশি থাকে না। % BN-SRC-016 TARGET-CORRECTION-BEGIN
\par\smallskip\noindent\textnormal{\small\textbf{উৎস-সংশোধন (BN-SRC-016):} হিমায়িত ইংরেজি উৎস বলে ব্যবধান প্রত্যেক ধাপে ঠিক অর্ধেক হয়। মধ্যবিন্দু নিজেই $S$-এর কোনো সদস্যের মান হলে উভয় প্রান্ত একই মধ্যবিন্দুতে চলে এসে ব্যবধান শূন্য হতে পারে; প্রমাণের জন্য যথেষ্ট এবং সর্বদা সত্য বক্তব্য হলো ব্যবধান সর্বাধিক অর্ধেক হয়।} % BN-SRC-016 TARGET-CORRECTION-END
তাই $\equivrep{f}{} = \equivrep{g}{}$।

প্রথমে দেখাই যে $\equivrep{f}{}$ হলো $S$-এর একটি ঊর্ধ্বসীমা, অর্থাৎ $(\forall h \in S)\equivrep{h}{} \leq \equivrep{f}{}$।
(এগোতে এগোতে আমরা \olref{thm:cauchyorderedfield} ব্যবহার করব।) ধরি $h \in S$ এবং স্ববিরোধে উপনীত হওয়ার জন্য ধরি $\equivrep{f}{} < \equivrep{h}{}$, ফলে $0_\Real < \equivrep{(h-f)}{}$। $f$ একটি একঘেয়ে হ্রাসমান কোশি অনুক্রম বলে এমন কোনো $n \in \Nat$ আছে যাতে $\equivrep{(c_{f(n)} - f)}{} < \equivrep{(h-f)}{}$। তাই:
\[
	(f(n))_\Real = \equivrep{c_{f(n)}}{} < \equivrep{f}{} + \equivrep{(h-f)}{} = \equivrep{h}{},
\]
কিন্তু নির্মাণ অনুসারে $\equivrep{h}{} \leq (f(n))_\Real$—এর সঙ্গে এটি বিরোধী।

এরপর দেখাই যে $\equivrep{f}{} = \equivrep{g}{}$ হলো $S$-এর \emph{লঘিষ্ঠ} ঊর্ধ্বসীমা। তাই $j$-কে যেকোনো কোশি অনুক্রম নিই এবং ধরি $\equivrep{j}{} < \equivrep{g}{}$। উপরের মতো যুক্তি দিয়ে ($g$ \emph{বর্ধমান}—এই তথ্য ব্যবহার করে) এমন $n \in \Nat$ আছে যাতে $\equivrep{j}{} < (g(n))_\Real$। কিন্তু নির্মাণ অনুসারে এমন $h \in S$ আছে যাতে $(g(n))_\Real \leq \equivrep{h}{}$; তাই $\equivrep{j}{} < \equivrep{h}{}$, এবং সেই কারণে $\equivrep{j}{}$ $S$-এর ঊর্ধ্বসীমা নয়।
\end{proof}

\end{document}
